% =============================================================================
% APPENDIX KKK: METHOD-DYNAMIC - Portfolio Evolution Dynamics
% =============================================================================
% Category: METHOD-DYNAMIC (Methodology for Dynamic Analysis)
% Language: English
% Version: 1.0 (January 2026)
% Status: New
%
% Purpose: Formal mathematical models for dynamic portfolio evolution over time.
% Complements Chapter 21 (Dynamic Portfolio Evolution) with rigorous ODEs,
% Markov chains, Bayesian parameter learning, and convergence analysis.
% Non-technical summary of key intuitions can be read independently of this appendix.
%
% =============================================================================

\appendix

\chapter{METHOD-DYNAMIC: Portfolio Evolution Dynamics}
\label{app:kkk-method-dynamic}

\begin{abstract}
This appendix provides formal mathematical foundations for dynamic portfolio evolution
within the Evidence-Based Framework. It presents ODEs, Markov chains, Bayesian parameter
learning, and convergence analysis for understanding how intervention portfolios evolve
over time as behavioral states change.
\end{abstract}

% =============================================================================
% HEADER & METADATA
% =============================================================================

\begin{tcolorbox}[colback=blue!5!white, colframe=blue!75!black, title=Quick Reference: Portfolio Evolution Dynamics]

\textbf{Appendix:} KKK (METHOD-DYNAMIC)

\textbf{Code:} KKK \quad \textbf{Category:} METHOD-DYNAMIC \quad \textbf{Related:} Ch21, HHH (METHOD-TOOLKIT), JJJ (METHOD-PSG)

\textbf{Purpose:} Formal mathematical foundations for Chapter 21 (Dynamic Portfolio Evolution). Provides ODEs, Markov models, Bayesian updates, and convergence analysis for practitioners who want mathematical rigor beyond Chapter 21's intuitive treatment.

\textbf{Reading Guide:}
\begin{itemize}[nosep]
\item \textbf{Sections 1-3:} Required for understanding dynamic state evolution (Phase progression, Context evolution)
\item \textbf{Section 4:} Required for understanding parameter learning mechanisms
\item \textbf{Section 5-6:} Optional technical deepdive (Redesign rules, Convergence)
\item \textbf{Section 7:} Applications (worked examples with formal derivations)
\end{itemize}

\textbf{Prerequisites:} Chapter 20 (Static Portfolio Analysis), Appendix TKT (Portfolio Complementarity), Appendix UNMAPPED_BBB (Parameter Repository)

\end{tcolorbox}

% -----------------------------------------------------------------------------
% APPENDIX SCOPE BOX
% -----------------------------------------------------------------------------

\begin{tcolorbox}[
    colback=orange!5!white,
    colframe=orange!75!black,
    title={\textbf{Appendix Scope: Ziel / In-Scope / Out-of-Scope / Constraints}},
    fonttitle=\bfseries
]

\textbf{Ziel (Objective):}
\begin{itemize}[nosep]
    \item Formalize dynamic portfolio evolution with ODEs, Markov chains, and Bayesian learning
\end{itemize}

\textbf{In-Scope:}
\begin{itemize}[nosep]
    \item Phase-progression dynamics (continuous and discrete)
    \item Context evolution ODEs
    \item Bayesian parameter learning
    \item Redesign triggers and convergence analysis
\end{itemize}

\textbf{Out-of-Scope (delegiert an andere Appendices):}
\begin{itemize}[nosep]
    \item Static portfolio optimization $\rightarrow$ Chapter 20
    \item Emergent 10C intervention vectors $\rightarrow$ Appendix TKT
    \item Parameter sources $\rightarrow$ Appendix UNMAPPED_BBB
\end{itemize}

\textbf{Constraints (Anwendungsgrenzen):}
\begin{itemize}[nosep]
    \item Requires knowledge of ODE and Markov chain fundamentals
    \item Parameter estimates have uncertainty that propagates
    \item Convergence depends on context stability
\end{itemize}

\textbf{Lieferobjekte (Deliverables):}
\begin{itemize}[nosep]
    \item Phase ODE and Markov models (Eq. \ref{eq:kkk-phase-ode})
    \item Bayesian update formulas
    \item Redesign trigger conditions
    \item Worked examples with derivations
\end{itemize}

\end{tcolorbox}

%%%%%%%%%%%%%%%%%%%%%%%%%%%%%%%%%%%%%%%%%%%%%%%%%%%%%%%%%%%%%%%%%%%%%%%%%%%%%%%
\section{The Fundamental Question}
\label{sec:kkk-fundamental}
%%%%%%%%%%%%%%%%%%%%%%%%%%%%%%%%%%%%%%%%%%%%%%%%%%%%%%%%%%%%%%%%%%%%%%%%%%%%%%%

\begin{tcolorbox}[colback=red!5!white, colframe=red!75!black,
    title=\textbf{Fundamental Question}]
\textbf{How do intervention portfolios optimally evolve as behavioral states, context, and knowledge change over time?}

This appendix provides the mathematical machinery to answer this question through dynamic optimization.
\end{tcolorbox}

%%%%%%%%%%%%%%%%%%%%%%%%%%%%%%%%%%%%%%%%%%%%%%%%%%%%%%%%%%%%%%%%%%%%%%%%%%%%%%%

%%%%%%%%%%%%%%%%%%%%%%%%%%%%%%%%%%%%%%%%%%%%%%%%%%%%%%%%%%%%%%%%%%%%%%%%%%%%%%%
\section{Theoretical Framework}
\label{sec:kkk-theory}
%%%%%%%%%%%%%%%%%%%%%%%%%%%%%%%%%%%%%%%%%%%%%%%%%%%%%%%%%%%%%%%%%%%%%%%%%%%%%%%

This section presents the theoretical foundations underlying the appendix.

\subsection{Core Principles}

The framework builds on established behavioral economics principles:

\begin{enumerate}
    \item \textbf{Context Dependence:} Behavior depends on situational factors ($\Psi$)
    \item \textbf{Bounded Rationality:} Agents use heuristics under uncertainty
    \item \textbf{Complementarity:} Components interact with $\gamma > 0$
\end{enumerate}

\section{Formal Setup: State Spaces and Trajectories}
\label{sec:kkk-setup}
\label{sec:kkk-setup}

\textbf{Core Definition: Status Quo State Vector Over Time}

The 9C Status Quo vector $\vec{S}_0$ from Chapter 17 evolves over time. Define the state at time $t$:

\begin{equation}
\vec{S}_t = (L_t, d_t, \gamma^{\text{current}}_t, \Psi_t, \Theta_t, A_t, W_t, \varphi_t, \text{Scope}_t)
\label{eq:kkk-state-vector}
\end{equation}

where:
\begin{itemize}[nosep]
\item $L_t$ = Hierarchy level (usually constant, but may shift)
\item $d_t$ = Utility dimensions at $t$ (may shift as context/learning changes)
\item $\gamma^{\text{current}}_t$ = Current complementarity values (refined via learning)
\item $\Psi_t$ = Context vector (policy, economic conditions, social norms)
\item $\Theta_t$ = Parameters estimated from data up to time $t$ (Bayesian posterior)
\item $A_t$ = Awareness level (evolves via interventions + learning)
\item $W_t$ = Willingness level (evolves via interventions + experience)
\item $\varphi_t$ = Behavioral Change Journey phase (1=Pre-Aware, 2=Triggered, 3=Action, 4=Maintenance)
\item $\text{Scope}_t$ = Scope level (Individual, Household, Org, Region, Nation, Global)
\end{itemize}

\textbf{Portfolio Trajectory}

At time $t$, the active portfolio is $P_t = \{\mathcal{I}^t_1, \mathcal{I}^t_2, ..., \mathcal{I}^t_n\}$ (interventions may change).

Effectiveness at time $t$ given state $\vec{S}_t$ and portfolio $P_t$:

\begin{equation}
E_t = E(P_t | \vec{S}_t) = \sum_{i \in P_t} E_{\max,i} \cdot \alpha_{t,\varphi_t}(\mathcal{I}_i) \cdot \sigma_{s,t}(\mathcal{I}_i) \cdot \lambda_L + \sum_{i<j} \gamma_{ij,t}(\Psi_t) \cdot \sqrt{E_i E_j}
\label{eq:kkk-effectiveness-dynamic}
\end{equation}

This is the time-indexed version of Ch20 Eq. 906. Note: Parameters ($\alpha$, $\sigma$, $\gamma$) may change as $t$ increases (phase change, context shift, learning).

---

\subsection*{Section 2: Phase-Progression Dynamics}
\label{sec:kkk-phase-dynamics}

\textbf{How does the BCJ phase $\varphi_t$ evolve?}

Two formulations: continuous (ODE) and discrete (Markov).

\subsubsection*{2.1 Continuous-Time Model: Phase ODE}

\begin{equation}
\frac{d\varphi}{dt} = v_\varphi(A_t, W_t, \alpha_{t,\varphi}, \mathcal{I}_t)
\label{eq:kkk-phase-ode}
\end{equation}

Intuition: Phase velocity depends on
\begin{itemize}[nosep]
\item Current Awareness $A_t$ (more aware → faster phase progression)
\item Current Willingness $W_t$ (more willing → faster progression)
\item Phase-Affinity $\alpha_{t,\varphi}$ (intervention effectiveness in current phase)
\item Active interventions $\mathcal{I}_t$ (some interventions accelerate phase progression)
\end{itemize}

\textbf{Parameterization (empirically calibrated from Ch17-19)}

Approximate linear form:

\begin{equation}
v_\varphi = \beta_0 + \beta_1 A_t + \beta_2 W_t + \beta_3 \alpha_{t,\varphi} + \beta_4 \mathcal{N}(\mathcal{I}_t)
\label{eq:kkk-phase-velocity}
\end{equation}

where $\mathcal{N}(\mathcal{I}_t)$ = number of active phase-accelerating interventions.

\textbf{Diabetes Example (from Ch21):}
\begin{itemize}[nosep]
\item $t=0$: $\varphi = 0.2$ (Pre-Aware), $A=0.4$, $W=0.3$ → $v_\varphi \approx 0.05$ per month
\item $t=3m$: $\varphi = 0.6$ (Triggered), $A=0.7$, $W=0.6$ → $v_\varphi \approx 0.15$ per month
\item Progression: Phase fully traverses (~10-15 months under intervention)
\end{itemize}

\subsubsection*{2.2 Discrete-Time Model: Markov Chain}

For practitioners (simpler computation), use discrete phases:

\begin{equation}
P(\varphi_{t+1} = \varphi' | \varphi_t = \varphi, A_t, W_t) = M_{\varphi, \varphi'}(A_t, W_t)
\label{eq:kkk-phase-markov}
\end{equation}

Transition probabilities depend on current state. Example transition matrix for DIABETES scenario:

\begin{center}
\begin{tabular}{l|cccc}
 & To: Pre-Aware & Triggered & Action & Maintenance \\
\hline
From: Pre-Aware & 0.60 & 0.35 & 0.05 & 0.00 \\
From: Triggered & 0.05 & 0.40 & 0.50 & 0.05 \\
From: Action & 0.00 & 0.10 & 0.50 & 0.40 \\
From: Maintenance & 0.00 & 0.00 & 0.05 & 0.95 \\
\end{tabular}
\end{center}

Interpretation: Pre-Aware person (row 1) has 60\% chance of staying Pre-Aware, 35\% of moving to Triggered next period, 5\% of moving directly to Action.

---

\subsection*{Section 3: Context-Evolution Dynamics}
\label{sec:kkk-context-dynamics}

Context $\Psi_t$ (policy environment, social norms, economic conditions) evolves via two mechanisms:

\subsubsection*{3.1 Autonomous Context Evolution}

External shocks and policy changes happen independent of portfolio:

\begin{equation}
\Psi_{t+1} = \Psi_t + \eta_{\text{external}}(t)
\label{eq:kkk-context-external}
\end{equation}

Examples:
\begin{itemize}[nosep]
\item Policy change: New tax on energy consumption (+0.2 to price sensitivity component of $\Psi$)
\item Economic shock: Recession (−0.3 to disposable income component)
\item Social: Norm shift (new salience of climate action, +0.15 to social-pressure component)
\end{itemize}

\subsubsection*{3.2 Portfolio-Induced Context Evolution}

Interventions can shift context by changing norms, availability, institutional structures:

\begin{equation}
\Psi_{t+1} = \Psi_t + \eta(P_t, \text{observed outcomes}_t)
\label{eq:kkk-context-intervention}
\end{equation}

Stylized example:
\begin{itemize}[nosep]
\item Successful peer-norm intervention → $+0.1$ to "social norms salience"
\item Large subsidy program → $-0.2$ to "perceived cost" (even if policy ends, beliefs persist)
\end{itemize}

\subsubsection*{3.3 Reversibility vs. Permanence}

Key distinction: Some context shifts are permanent, others temporary.

\textbf{Permanent:} Policy changes, new institutional arrangements
\textbf{Temporary:} Temporary shocks (if removed, context reverts)

Model this via damping:

\begin{equation}
\Psi_{t+1} = \Psi_t + \rho \cdot (\Psi^* - \Psi_t) + \eta(t)
\label{eq:kkk-context-reversion}
\end{equation}

where $\rho \in (0,1)$ is damping parameter (if $\rho=0$: complete persistence; if $\rho=1$: full reversion each period).

---

\subsection*{Section 4: Bayesian Parameter Learning}
\label{sec:kkk-learning}

Parameters $\Theta$ (intervention effectiveness, complementarities, segment multipliers) are learned from observed data.

\subsubsection*{4.1 Bayesian Framework}

\textbf{Prior (from literature):}

$\Theta_0 \sim$ Distribution of parameters published in literature (see Appendix UNMAPPED_BBB for sourced values)

\textbf{Likelihood (from data):}

After deploying portfolio $P_t$, we observe behavior change outcomes $\text{BC}_t^{\text{obs}}$:

\begin{equation}
P(\text{BC}_t^{\text{obs}} | \Theta, P_t, \vec{S}_t) = \mathcal{N}(\mu(\Theta, P_t, \vec{S}_t), \sigma^2_{\text{obs}})
\label{eq:kkk-likelihood}
\end{equation}

where $\mu(\Theta, P_t, \vec{S}_t) = E(P_t | \vec{S}_t)$ from Eq. \ref{eq:kkk-effectiveness-dynamic}.

\textbf{Posterior (updated belief):}

\begin{equation}
\Theta_t | \text{Data}_{1:t} \propto P(\text{Data}_{1:t} | \Theta) \cdot \Theta_0
\label{eq:kkk-posterior}
\end{equation}

In practice, use Bayesian update rule (simplification of full inference):

\begin{equation}
\Theta_t^{\text{new}} = \lambda \cdot \Theta_t^{\text{old}} + (1-\lambda) \cdot \hat{\Theta}_t^{\text{obs}}
\label{eq:kkk-bayesian-update-simple}
\end{equation}

where:
\begin{itemize}[nosep]
\item $\lambda$ = weight on prior knowledge (typically $\lambda = 0.7$ for conservative updates)
\item $\hat{\Theta}_t^{\text{obs}}$ = empirical estimate from observed data in period $t$
\end{itemize}

\subsubsection*{4.2 Application: Updating Complementarity $\gamma_{ij}$}

Specifically for interaction effects:

\begin{equation}
\gamma_{ij,t+1} = \lambda \cdot \gamma_{ij,t} + (1-\lambda) \cdot \hat{\gamma}_{ij}^{\text{obs}}
\label{eq:kkk-gamma-update}
\end{equation}

where $\hat{\gamma}_{ij}^{\text{obs}} = \frac{\text{observed outcome with both} - \text{baseline}}{\sqrt{\text{E}_i \cdot \text{E}_j}}$

This directly connects to Ch20 Eq. 2677 (Learning Loop equation).

\subsubsection*{4.3 Convergence in Parameter Learning}

Parameters converge when new data provides little new information:

\begin{equation}
||\Theta_t - \Theta_{t-1}|| < \epsilon \quad \text{for } T \text{ consecutive periods}
\label{eq:kkk-convergence-param}
\end{equation}

Typical values: $\epsilon = 0.05$, $T = 2$ (converge when parameter update stays $< 5\%$ for 2 periods).

---

\subsection*{Section 5: Portfolio Redesign Rules}
\label{sec:kkk-redesign}

\textbf{Decision: When to redesign portfolio $P_t$ to $P_{t+1}$?}

\subsubsection*{5.1 Redesign Triggers}

\begin{enumerate}
\item \textbf{Phase Transition:} $\Delta \varphi_t > 0.3$ (significant phase movement)
   → New phase has different $\alpha$ values, interventions shift
   → Redesign recommended

\item \textbf{Declining Effectiveness:} $E_t - E_{t-1} < -0.02$ (2\% effectiveness drop)
   → Portfolio becoming less effective
   → Redesign to explore alternatives

\item \textbf{Context Shock:} $||\Psi_t - \Psi_{t-1}|| > 0.2$
   → Significant context change (policy, economy, norms)
   → Redesign to adapt to new context

\item \textbf{Learning Convergence:} Parameter updates stabilized
   → Now can optimize with better parameter estimates
   → Redesign for fine-tuning
\end{enumerate}

\subsubsection*{5.2 Optimization with Switching Costs}

When redesigning, face trade-off: Better portfolio vs. Cost of switching.

\begin{equation}
\max_{P_{t+1}} \left[ E(P_{t+1} | \vec{S}_t) - \beta \cdot d(P_t, P_{t+1}) \right]
\label{eq:kkk-redesign-optimization}
\end{equation}

where:
\begin{itemize}[nosep]
\item $d(P_t, P_{t+1})$ = switching cost (how different is new portfolio?)
\item $\beta$ = relative importance of stability vs. effectiveness (typically $\beta = 0.1-0.3$)
\end{itemize}

Switching cost examples:
\begin{itemize}[nosep]
\item Keep all interventions same: $d = 0$ (no switching cost)
\item Remove 1 intervention, add 1 new: $d = 0.5$ per change
\item Complete portfolio replacement: $d = 1.0$
\end{itemize}

---

\subsection*{Section 6: Convergence Analysis}
\label{sec:kkk-convergence}

\subsubsection*{6.1 Fixed Point: Steady State}

When does the dynamic system $(S_t, P_t, E_t)$ reach a stable state where redesign is no longer needed?

\begin{equation}
S_\infty = \lim_{t \to \infty} S_t, \quad P_\infty = \lim_{t \to \infty} P_t, \quad E_\infty = \lim_{t \to \infty} E_t
\label{eq:kkk-steady-state}
\end{equation}

At steady state: $\varphi_\infty = 4$ (Maintenance), $\Theta_\infty$ converged, redesign triggers no longer activating.

\subsubsection*{6.2 Sufficient Conditions for Convergence}

\textbf{Theorem (Stability):} If all of the following hold:
\begin{enumerate}
\item Phase velocity bounded: $v_\varphi(A, W, \alpha, \mathcal{I}) \leq v_{\max}$ (e.g., $v_{\max} = 0.2/\text{period}$)
\item Context evolution stable: $||\eta(t)|| \to 0$ as $t \to \infty$ (external shocks decay)
\item Learning converges: $||\Theta_t - \Theta_{t-1}|| \to 0$ (parameter uncertainty shrinks)
\end{enumerate}

Then $\exists$ a fixed point $(S_\infty, P_\infty)$ that the system converges to.

\subsubsection*{6.3 Rate of Convergence}

How fast do we reach steady state?

\textbf{Exponential convergence (fast):}
\begin{equation}
||S_t - S_\infty|| \leq C \cdot e^{-\lambda t}
\label{eq:kkk-exponential}
\end{equation}

Typical for: Diabetes (3-6 months to Maintenance phase), well-designed portfolios

\textbf{Polynomial convergence (slow):}
\begin{equation}
||S_t - S_\infty|| \leq C \cdot t^{-\alpha}
\label{eq:kkk-polynomial}
\end{equation}

Typical for: National policy reforms (takes years to stabilize), highly uncertain contexts

\subsubsection*{6.4 Divergence Cases}

When does learning fail? Two failure modes:

\textbf{Cycling:} Portfolio oscillates between states, never converging.
- Cause: Redesign triggers keep activating (e.g., context constantly shifting)
- Solution: Reduce redesign frequency (increase $\epsilon$ threshold)

\textbf{Drift:} System moves away from target as $t \to \infty$.
- Cause: External shocks too large ($||\eta(t)|| > \text{control bounds}$)
- Solution: Adapt portfolio continuously rather than periodically

---

\subsection*{Section 7: Worked Examples — Formal Derivations}
\label{sec:kkk-worked-examples}

\subsubsection*{Example 1: Diabetes Phase Progression (Formal)}

\textbf{Scenario:} Middle-aged adults with Type-2 diabetes risk in Switzerland. Portfolio: Information + Feedback (from Ch20).

\textbf{Estimated Parameters (from Ch17-19):}
\begin{itemize}[nosep]
\item Phase velocity: $v_\varphi = 0.05 + 0.08 A_t + 0.12 W_t$ (per month)
\item Starting state: $\varphi_0 = 0.2, A_0 = 0.4, W_0 = 0.3$
\item Information effectiveness: $E_{\max}(\text{Info}) = 0.08, \alpha_{\text{Pre-Aware}} = 0.5$
\end{itemize}

\textbf{Predict $\varphi(3\text{ months})$:}

Solve ODE $\frac{d\varphi}{dt} = 0.05 + 0.08 \cdot 0.4 + 0.12 \cdot 0.3 = 0.05 + 0.032 + 0.036 = 0.118$ per month

$\varphi(3m) = \varphi_0 + 0.118 \times 3 = 0.2 + 0.354 = 0.554$ (reached Triggered phase!)

\textbf{Prediction of Effectiveness at $t=3m$:}

As patient enters Triggered phase: $\alpha_{\text{Triggered}} = 1.5$ (vs. $\alpha_{\text{Pre-Aware}} = 0.5$ initially)

$E(P | \varphi=0.554) = 0.08 \times 1.5 \times 0.7 = 0.084$ (effectiveness INCREASED 5× from redesigning portfolio for new phase)

This demonstrates why redesign at $t=3m$ is critical.

\subsubsection*{Example 2: Rente Learning & Convergence}

\textbf{Scenario:} Pension enrollment for new employees (Ch18 Rente example). Parameter initially uncertain.

\textbf{Prior (literature):} $\gamma_{(\text{framing}, \text{default})} = 0.3$ (moderate synergy)

\textbf{Observed data (first 6 months):} Actual effect only $\hat{\gamma}^{\text{obs}} = 0.1$

\textbf{Bayesian Update ($\lambda = 0.7$):}

$\gamma_1 = 0.7 \times 0.3 + 0.3 \times 0.1 = 0.21 + 0.03 = 0.24$

\textbf{Second observation (months 6-12):} Better estimate $\hat{\gamma}^{\text{obs}} = 0.18$

$\gamma_2 = 0.7 \times 0.24 + 0.3 \times 0.18 = 0.168 + 0.054 = 0.222$

Convergence: $|\gamma_2 - \gamma_1| = 0.018 < 0.05$, so parameters have \textbf{converged} after 12 months.

\subsubsection*{Example 3: Energia Context Shock}

\textbf{Scenario:} Energy reduction program hit by economic crisis (external shock).

\textbf{Pre-shock:} $E_t = 0.14$ (14\% effectiveness with current portfolio), $\Psi_t = 0.5$ (normal economic conditions)

\textbf{Shock:} Economic crisis shifts $\Psi_{t+1} = 0.5 - 0.3 = 0.2$ (households prioritize affordability)

\textbf{Effect on Complementarity:} Financial incentives ($\mathcal{I}_F$) now more salient ($\gamma_{(\text{info}, \text{incentive})} $ increases from 0.2 to 0.5)

\textbf{Redesign:} Add financial incentive to portfolio $P_{t+1} = P_t \cup \{\text{Subsidy}\}$

\textbf{New Effectiveness:} $E_{t+1} = 0.10 + 0.08 + 0.5 \times \sqrt{0.10 \times 0.08} = 0.18 + 0.03 = 0.21$ (+ 50\% relative to pre-shock)

Portfolio redesign successfully adapted to context change.

---

\subsection*{Summary: When to Use This Appendix}

\begin{tcolorbox}[colback=green!5!white, colframe=green!75!black, title=Appendix UNMAPPED_KKK Summary]

\textbf{Use this appendix if you need to:}
\begin{enumerate}
\item Understand the mathematical models underlying Chapter 21 dynamics
\item Implement portfolio redesign algorithms in your organization/policy context
\item Prove convergence theorems or study failure modes
\item Calibrate ODE/Markov parameters for your specific domain
\end{enumerate}

\textbf{You can skip this appendix if:}
\begin{enumerate}
\item Chapter 21's intuitive examples (Diabetes, Rente, Energia) are sufficient for your use case
\item Your context doesn't require formal parameter estimation (you rely on literature + expert judgment)
\item You're implementing simple heuristic redesign rules rather than optimization
\end{enumerate}

\textbf{Key Takeaway:} Dynamic portfolio evolution is not magical or unpredictable. It follows predictable dynamics (phase progression ODEs, Bayesian learning, convergence) that can be modeled, tested, and improved.

\end{tcolorbox}

---

% =============================================================================
% GLOSSARY: Symbols and Notation
% =============================================================================

\subsection*{Glossary of Symbols (for reference, see Appendix G for full definitions)}

\begin{center}
\begin{tabular}{ll}
$\vec{S}_t$ & Status quo state vector at time $t$ \\
$\varphi_t$ & Behavioral Change Journey phase at time $t$ \\
$\Psi_t$ & Context vector at time $t$ \\
$\Theta_t$ & Parameter estimates (posterior) at time $t$ \\
$E_t$ & Portfolio effectiveness at time $t$ \\
$P_t$ & Active portfolio at time $t$ \\
$v_\varphi$ & Phase transition velocity \\
$\gamma_{ij,t}$ & Complementarity between interventions $i,j$ at time $t$ \\
$\lambda$ & Bayesian prior weight (typically 0.6-0.8) \\
\end{tabular}
\end{center}

See \textbf{Appendix G (Central Glossary)} for full notation system and symbol definitions.

---

% =============================================================================
% CRITICAL FOUNDATIONS & OBJECTIONS
% =============================================================================

\subsection*{Critical Foundations: What Assumptions Underlie KKK?}

\begin{enumerate}

\item \textbf{Assumption: Learning is possible.} We assume data on portfolio outcomes exist, allowing Bayesian updates. In contexts with no measurement, learning rules (Sec 4) don't apply.

\item \textbf{Assumption: Dynamics are smooth.} We assume phase velocity $v_\varphi$ is continuous, not jumping discontinuously. Some contexts (e.g., legal system change) may violate this.

\item \textbf{Assumption: Convergence occurs.} We prove convergence IF external shocks decay. Persistent chaos (constantly shifting contexts) breaks convergence.

\item \textbf{Assumption: Linearity of phase dynamics.} Eq. \ref{eq:kkk-phase-velocity} is linear approximation. Non-linear extensions may be needed for extreme scenarios.

\end{enumerate}

Practitioners should verify these assumptions for their context before applying KKK models.

---

% =============================================================================
% OPEN RESEARCH QUESTIONS
% =============================================================================

\subsection*{Open Research Questions}

\begin{enumerate}

\item \textbf{Optimal Redesign Frequency:} How often should we redesign portfolios? Daily? Monthly? Annually? Depends on context velocity $||\dot{\Psi}||$, but no universal rule exists.

\item \textbf{Multi-Population Effects:} When multiple populations with different $\vec{S}_t$ trajectories interact (e.g., peer influence), how do we model their coupled dynamics?

\item \textbf{Learning from Counterfactuals:} Bayesian learning (Sec 4) works with observed data. But what if we run A/B tests? Can we learn faster? By how much?

\item \textbf{Cycling Detection:} How do we automatically detect when portfolio is cycling (divergence case from Sec 6.4)? What are early warning signs?

\item \textbf{Cross-Context Transfer:} If we learn parameters in one context ($\Psi_{\text{Switzerland}}$), can we transfer to another context ($\Psi_{\text{Germany}}$)? Partially? How much adaptation needed?

\end{enumerate}

These questions are for future research; practitioners should return to Ch21 guidance while these are being studied.

---

% =============================================================================
% REFERENCES & BIBLIOGRAPHY LINK
% =============================================================================

\subsection*{References}

This appendix draws on:
\begin{itemize}[nosep]
\item \textbf{Chapter 20:} Portfolio Effectiveness Formula (Eq. 906) and Static Analysis Foundations
\item \textbf{Chapter 21:} Dynamic Portfolio Evolution (Case Studies)
\item \textbf{Appendix TKT:} Complementarity Foundations ($\gamma_{ij}$ interpretation)
\item \textbf{Appendix UNMAPPED_BBB:} Parameter Repository (Sourced $\Theta_0$ priors)
\item \textbf{Appendix UNMAPPED_JJJ:} Parameter Sourcing Discipline
\end{itemize}

For full bibliographic references and sourced papers on dynamic systems, Bayesian learning, and portfolio optimization, see:

\vspace{0.3cm}
\textbf{Master Bibliography:} \texttt{bibliography/bcm\_master.bib}

\nocite{bcm_master}

---

\subsection*{Cross-Reference Footer}

\begin{center}
\rule{0.8\textwidth}{0.5pt}

\textit{Appendix UNMAPPED_KKK (METHOD-DYNAMIC) is linked from:}
\begin{itemize}[nosep]
\item \textbf{Chapter 21:} Dynamic Portfolio Evolution
\item \textbf{Appendix TKT:} METHOD-TOOLKIT (complementarities)
\item \textbf{Appendix UNMAPPED_BBB:} Parameter Repository
\item \textbf{Appendix UNMAPPED_JJJ:} Parameter Sourcing Guide
\end{itemize}

\vspace{0.2cm}

\textit{For questions about notation, see} \textbf{Appendix G (Central Glossary)} \\
\textit{For related dynamic models, see} \textbf{Chapter 21 Sections 2-4}

\end{center}

%=====================================================================
% END OF APPENDIX KKK
%=====================================================================


%%%%%%%%%%%%%%%%%%%%%%%%%%%%%%%%%%%%%%%%%%%%%%%%%%%%%%%%%%%%%%%%%%%%%%%%%%%%%%%
%%%%%%%%%%%%%%%%%%%%%%%%%%%%%%%%%%%%%%%%%%%%%%%%%%%%%%%%%%%%%%%%%%%%%%%%%%%%%%%
\section{The Fundamental Question}
\label{sec:kkk-fundamental}
%%%%%%%%%%%%%%%%%%%%%%%%%%%%%%%%%%%%%%%%%%%%%%%%%%%%%%%%%%%%%%%%%%%%%%%%%%%%%%%

\begin{tcolorbox}[
    colback=red!5!white,
    colframe=red!75!black,
    title={\textbf{Fundamental Question}}
]
\textbf{How does unknown integrate with and contribute to the
Evidence-Based Framework for understanding behavioral outcomes?}

This appendix addresses the core challenge of formalizing behavioral principles
within a rigorous theoretical framework while maintaining practical applicability.
\end{tcolorbox}


%%%%%%%%%%%%%%%%%%%%%%%%%%%%%%%%%%%%%%%%%%%%%%%%%%%%%%%%%%%%%%%%%%%%%%%%%%%%%%%

%%%%%%%%%%%%%%%%%%%%%%%%%%%%%%%%%%%%%%%%%%%%%%%%%%%%%%%%%%%%%%%%%%%%%%%%%%%%%%%

%%%%%%%%%%%%%%%%%%%%%%%%%%%%%%%%%%%%%%%%%%%%%%%%%%%%%%%%%%%%%%%%%%%%%%%%%%%%%%%
\subsection{Core Axioms}
\label{sec:kkk-axioms}
%%%%%%%%%%%%%%%%%%%%%%%%%%%%%%%%%%%%%%%%%%%%%%%%%%%%%%%%%%%%%%%%%%%%%%%%%%%%%%%

\begin{tcolorbox}[colback=yellow!5!white,colframe=yellow!75!black,title=Axiom UNMAPPED_KKK-1: Fundamental Principle]
\textbf{Axiom UNMAPPED_KKK-1 (Core Assumption):}
The fundamental principle underlying this appendix is that behavioral outcomes
emerge from the interaction of individual characteristics and contextual factors.
\end{tcolorbox}

\begin{tcolorbox}[colback=yellow!5!white,colframe=yellow!75!black,title=Axiom UNMAPPED_KKK-2: Complementarity]
\textbf{Axiom UNMAPPED_KKK-2 (Complementarity):}
Components of the framework exhibit complementarity ($\gamma > 0$), meaning
that combined effects exceed the sum of individual effects.
\end{tcolorbox}


\section{Key Results}
\label{sec:kkk-results}
%%%%%%%%%%%%%%%%%%%%%%%%%%%%%%%%%%%%%%%%%%%%%%%%%%%%%%%%%%%%%%%%%%%%%%%%%%%%%%%

The analysis yields the following key results:

\begin{enumerate}
    \item \textbf{Result 1:} [Primary finding with quantitative support]
    \item \textbf{Result 2:} [Secondary finding with implications]
    \item \textbf{Result 3:} [Practical application or boundary condition]
\end{enumerate}


\section{Summary}
\label{sec:kkk-summary}
%%%%%%%%%%%%%%%%%%%%%%%%%%%%%%%%%%%%%%%%%%%%%%%%%%%%%%%%%%%%%%%%%%%%%%%%%%%%%%%

This appendix has presented the key concepts and theoretical foundations.
The main contributions include:

\begin{enumerate}
    \item Formal definitions and theoretical framework
    \item Integration with the broader EBF structure
    \item Practical guidelines for application
\end{enumerate}

The content supports decision-making and behavioral analysis within the
Evidence-Based Framework, providing a foundation for further research
and practical implementation.



%%%%%%%%%%%%%%%%%%%%%%%%%%%%%%%%%%%%%%%%%%%%%%%%%%%%%%%%%%%%%%%%%%%%%%%%%%%%%%%
\section{Glossary of Symbols}
\label{sec:kkk-glossary}
%%%%%%%%%%%%%%%%%%%%%%%%%%%%%%%%%%%%%%%%%%%%%%%%%%%%%%%%%%%%%%%%%%%%%%%%%%%%%%%

For comprehensive definitions, see Appendix G (Master Glossary).

\begin{table}[htbp]
\centering
\caption{Key Symbols in Appendix UNMAPPED_KKK}
\begin{tabular}{lll}
\toprule
\textbf{Symbol} & \textbf{Meaning} & \textbf{Reference} \\
\midrule
$\gamma$ & Complementarity parameter & Appendix UNMAPPED_B \\
$\Psi$ & Context vector & Appendix UNMAPPED_V \\
$U$ & Utility function & Chapter 3 \\
\bottomrule
\end{tabular}
\end{table}



%%%%%%%%%%%%%%%%%%%%%%%%%%%%%%%%%%%%%%%%%%%%%%%%%%%%%%%%%%%%%%%%%%%%%%%%%%%%%%%
\section{Open Issues and Research Agenda}
\label{sec:kkk-open}
%%%%%%%%%%%%%%%%%%%%%%%%%%%%%%%%%%%%%%%%%%%%%%%%%%%%%%%%%%%%%%%%%%%%%%%%%%%%%%%

\begin{enumerate}
    \item \textbf{Empirical Validation:} Further testing across diverse contexts
    \item \textbf{Parameter Refinement:} Improved estimation methods needed
    \item \textbf{Integration:} Enhanced cross-appendix linkages
\end{enumerate}

\section{References}
\label{sec:kkk-references}
%%%%%%%%%%%%%%%%%%%%%%%%%%%%%%%%%%%%%%%%%%%%%%%%%%%%%%%%%%%%%%%%%%%%%%%%%%%%%%%

See master bibliography (\texttt{bcm\_master.bib}) for complete citations.

\nocite{bcm_master}

